% 第 31 章　统计物理：用概率描述大军团
\chapter{统计物理：用概率描述大军团}

\step{反常识开场}

给自行车打气，胎压表指针停在一个刻度上，纹丝不动。这看起来再平常不过，可只要你多想一秒，它就是一个奇迹。

轮胎里的空气不是一堵安静的墙。按照热学的分子图像，它是几千亿亿个飞来飞去的小球，每个都以每秒四五百米的速度狂奔，不断撞上轮胎内壁再弹开。压强，就是这些撞击的总效果。我们来数一数：在通常的气压下，每平方厘米的器壁，每秒要承受大约 \(10^{23}\) 次撞击。这个数写出来是一后面跟着二十三个零。每一次撞击的时机、位置、角度都毫无规律可言，没有任何一个分子知道自己“应该”撞哪里。

那么问题来了：一支由完全无规律的个体组成的军团，为什么能给出稳定到连精密仪器都看不出波动的总效果？指针为什么不抖？把场面再放大一点：整个大气层压在每平方厘米地面上的，也是这样一支乱军的总账，可你从未觉得耳膜在被随机地捶打。更进一步：混乱怎么能“制造”出秩序，而且是最可靠、最铁面无私的那种秩序？

这不是哲学问题，而是一条可以用硬币、骰子和一杯空气亲自验证的物理规律。其实它的影子你早就见过：一座千万人口的城市，没人指挥，每天的用电量、用水量、急诊人数却稳定得可以提前一周做计划；而同一个办公室只有十个人时，谁请假、谁加班就完全没法预料。个体的不可预测与整体的稳定并存，这是同一个事实的两面。这一章我们就要学会一种新的思考方式：不追踪任何一个粒子，只数状态，照样能预言压强、温度、化学反应的快慢，甚至解释冰为什么融化、太阳为什么烧得慢。

\step{先猜一猜}

在翻开答案之前，请把你的直觉写下来，本章结尾我们再对照。

第一猜：把十枚硬币同时抛一次，正好五正五反的可能性有多大？很多人的第一反应是“一半对一半最平均，概率应该过半”，也有人猜三分之一。

第二猜：把一千枚硬币同时抛一次，正面的比例落在 \(49\%\) 到 \(51\%\) 之外的可能性有多大？直觉会说“数量越多越接近一半一半，超出这个范围应该很罕见”。

第三猜：你现在待的房间里，空气分子会不会在某个瞬间全部挤到房间的左半边，让你右半边的鼻子吸不到空气？注意，这件事不违反任何力学定律——每个分子往左飞并不犯法。大多数人回答“永远不会”。可如果每个分子都随机乱飞，凭什么“永远不会”？你的底气从哪里来？

这三道题的答案分别藏在概率分布、涨落和微观态计数里。写好了吗？写好再往下读——直觉只有先押了注，错了才会有痛感，而痛感是最好的老师。我们先做实验。

\step{动手实验}

\begin{experiment}[硬币大军]
\textbf{材料：}十枚硬币（或十枚同样的纽扣、瓶盖）、纸和笔。

\textbf{步骤：}
\begin{enumerate}[itemsep=2pt]
  \item 把十枚硬币握在手里摇一摇，撒在桌上，数一数正面朝上的个数，记在纸上。这算一次。
  \item 重复五十次。每次都记录正面数，会得到五十个 \(0\) 到 \(10\) 之间的数字。
  \item 画直方图：横轴标出 \(0,1,2,\dots,10\)，统计每个正面数出现了几回，画成柱子。
  \item 把这五十个数字加起来除以五十，得到平均值。
  \item 加做一步：把五十次按顺序分成五组，每组十次，分别算每组的平均值，看看这五个平均值一不一样。
\end{enumerate}

\textbf{你会看到什么：}单独看任何一次，数字乱跳，完全猜不到下一次是几；可五十次画成的直方图却有稳定的形状——中间高、两头低，大多数结果落在 \(3\) 到 \(7\) 之间，正好十枚正反对半分的情况大约只占四分之一，而“十枚全正面”你大概率一次也遇不到。五个小组的平均值也互不相同，在 \(5\) 附近上下浮动。如果和同桌的结果对比，你们两张直方图的细节处处不同，形状却像一个模子刻出来的——细节属于运气，形状属于规律。

\textbf{这说明什么：}你没有控制任何一枚硬币，却得到了一条可以预测、可以画在纸上的规律。规律不属于单个硬币，属于整个大军团。这就是统计物理的入口。
\end{experiment}

\step{旧模型遇到麻烦}

在本书的力学部分，我们习惯了一种世界观：只要知道一个物体此刻的位置和速度，以及它受到的力，就能算出它的未来。对一颗台球，这套方法好用到近乎完美。那为什么不对轮胎里的每个分子照做一遍？

第一道麻烦是数量。一升空气里大约有 \(2.7\times10^{22}\) 个分子。就算每个分子只记三个位置坐标和三个速度分量，也需要一台比全世界所有计算机加起来还大的机器才能存下这些数据，更别提测量它们了——任何测量本身就会打扰被测的分子。这个数字大到什么程度？一杯水里的分子数，比地球上所有海滩的沙子粒数还多出许多个数量级；把它们排成一列，可以绕地球到太阳之间往返无数次。

第二道麻烦更本质：多体系统对初始条件极端敏感。两个分子的速度差一点点，碰撞之后这点差别会被放大；经过几十次碰撞，误差就大到让预言完全失效。这在力学里叫混沌。也就是说，哪怕我们真的测得了某一瞬间的全部数据，详细轨迹的预言也会在极短时间内作废。

第三道麻烦其实是解放：我们根本不想要那么细的答案。轮胎厂关心的是压强，不是第 \(8000\) 亿亿号分子现在在哪。流行病学家关心的是全城发热人数的曲线，而不是张三今天会不会咳嗽。交通工程师关心的是某座桥每分钟通过多少辆车，而不是哪辆车几点几分上桥。对大军团提出单兵式的问题，本身就是问错了问题。

这里值得请出一个比喻：保险公司。它像本章的地方在于——保险公司不知道、也不需要知道哪个客户明年出车祸，但靠着大数规律，它能把明年的赔付总额算得相当准，照样赚钱。它不像本章的地方在于——人会看广告、会搬家、会改变习惯，个体行为会随着统计本身而变化；分子没有意图，它们的“统计规律”背后是严格的力学和能量守恒在约束，因而比任何社会统计都坚硬。比喻到此为止，接下来要用物理自己的语言。

还有一点必须分清：压强指针稳定，这是实验事实；“分子在做永不停歇的无规则运动”是我们对事实的物理解释；而“正面数的直方图呈某种数学形状”是数学模型。三者各有来源，谁也不能冒充谁。这一章的每一步都要记得自己站在哪一层。

\step{发明新概念}

现在像十九世纪的物理学家一样，动手发明新工具。我们从实验里的硬币开始，把一个例子向下挖五层。

\textbf{看见它：}直方图中间高两头低，这是形状，不是偶然。

\textbf{说清它：}关键变量是正面数。注意一个重要区分——“十枚硬币里正好五枚正面”这句话描述的是一个\textbf{宏观态}：它只说总数，不说哪几枚正面。而“第 \(1,3,4,7,9\) 号硬币正面，其余反面”描述的是一个\textbf{微观态}：每个个体的情况都指定了。一个宏观态可以由许许多多不同的微观态实现。五正五反这个宏观态，对应的微观态有 \(252\) 种排法；而十枚全正面，只有一种排法。

\textbf{画出它：}把 \(0\) 到 \(10\) 每个正面数对应的排法数列出来：\(1,10,45,120,210,252,210,120,45,10,1\)，再画成图，就解释了直方图的形状。

\begin{center}
\begin{tikzpicture}[x=0.95cm,y=0.85cm]
  \draw[-{Stealth}] (-0.4,0) -- (11.6,0);
  \node[below] at (5.5,-0.1) {正面数};
  \draw[-{Stealth}] (0,-0.3) -- (0,3.1) node[above] {排法数};
  \foreach \x/\h in {0/0.01,1/0.10,2/0.45,3/1.20,4/2.10,5/2.52,6/2.10,7/1.20,8/0.45,9/0.10,10/0.01}
    \fill[blue!45] (\x-0.32,0) rectangle (\x+0.32,\h);
  \foreach \x in {0,1,2,3,4,5,6,7,8,9,10} \node[below] at (\x,0) {\scriptsize \x};
  \node at (5,2.75) {\scriptsize \(252\)};
  \node at (4,2.32) {\scriptsize \(210\)};
  \node at (6,2.32) {\scriptsize \(210\)};
  \node[align=left] at (9.4,2.5) {\scriptsize 十枚硬币各正面数\\[-2pt]\scriptsize 对应的微观态数目};
\end{tikzpicture}
\end{center}

\textbf{算出它：}每种排法的机会均等（硬币没有偏好），所以某个正面数的概率就等于它的排法数除以总排法数 \(2^{10}=1024\)。这里藏着一个必须说破的假设：凭什么认定每种微观态机会均等？答案不是逻辑，而是对称性——硬币的两面没有任何区别，交换任意两枚硬币，物理规律一个字也不变。凡是找不到区别对待的理由，就平等对待：这是统计物理里最大胆也最好用的一步，它的正确性最终只能靠实验背书。五正五反的概率是 \(252/1024\approx 24.6\%\)——这就是第一猜的答案，远没过半。直觉错在哪？直觉把“最可能”当成了“很可能”。五正五反确实是最可能的单一结果，但它仍然要和其他九百多种排法分蛋糕。

\textbf{追到底：}把 \(10\) 枚换成 \(N\) 枚，同一条逻辑继续成立：各微观态机会均等，宏观态的概率正比于它包含的微观态数目 \(W\)。再换一个道具验证一遍：掷两颗骰子，点数和为 \(7\) 有 \(6\) 种掷法，和为 \(2\) 只有 \(1\) 种，所以赌场里“七”出现得最勤——古人没学过排列组合，却从无数次掷骰里摸出了这条计数规律。这就是整个统计物理的地基，玻尔兹曼把它刻在了自己的墓碑上：自然界不追踪任何一个粒子，它“选择”的永远是实现方式最多的那种宏观状态。房间里空气均匀分布，不是因为分子们爱好和平，而是因为“均匀”对应的微观态数目比“全挤在左边”多了天文数字倍——多到后者在宇宙年龄内都不会出现一次。

先把视线从硬币移开，看看这条逻辑管多宽。一个班同学的身高、一次考试全年级每一题的得分率、晚高峰路口每分钟通过的车辆数，画成直方图，常常都是中间高、两头低的钟形。为什么形状如此相似？因为这些量都是“许多个互不相干的小因素加起来的总和”：身高是遗传、营养、睡眠等一大堆小影响叠加的结果；路口车流量是成千上万个司机各自决定几点下班的叠加。数学里有一条中心极限定理保证：只要叠加的因素足够多、彼此不相干，总和的分布就趋向同一个钟形。这又是一个三层分明的例子——直方图的形状是实验事实，中心极限定理是数学模型，“身高由许多小因素叠加”是物理解释。钟形曲线无处不在，不是因为世界爱画钟，而是因为“加总”这个动作无处不在。

接下来发明第二个概念：\textbf{涨落}。抛十枚硬币，正面数常在 \(3\) 到 \(7\) 之间晃；抛一百万枚，正面数通常在四十九万九千到五十万一千之间晃。注意两件事：晃动的绝对范围变大了（从几枚变成上千枚），但晃动的相对范围急剧变小。一条普遍的经验规律是：绝对涨落大约按 \(\sqrt{N}\) 增长，相对涨落大约按 \(1/\sqrt{N}\) 缩小。这就是大数定律的定量面孔——大军团的稳定性不是来自个体变乖，而是来自偏离被人头摊薄。

第三个概念回答“为什么高能的分子少”。想象总能量固定的一盒气体，能量像一笔定额的奖金分给所有分子。把大笔奖金集中到少数分子头上的分法，种类很少；把奖金大致摊平的分法，种类极多。既然自然界按微观态数目“投票”，能量越高的分子自然越稀少。这个比喻像的地方在于：分法确实有穷尽和计数，总额确实守恒。它不像的地方在于：钱可以任意拆分，而分子的能量交换只能通过一次次碰撞进行，而且在量子世界里能量常常只能整份整份地换手——正是这个“整份”埋下了第三册量子论的伏笔。这个直觉马上要变成一条公式。

\step{一句话模型}

三个新概念已经到手：概率分布告诉我们每种结果占多大份额，涨落告诉我们平均值的权威随粒子数按 \(1/\sqrt{N}\) 增长，微观态计数告诉我们自然界偏爱哪种宏观状态。三者拧成一股绳，就是本章的核心，可以压缩成一句话：

\begin{oneline}
单个粒子怎么走，我们管不了、也不必管；大量粒子的集体行为由“哪种宏观状态对应的微观实现方式最多”决定。数状态数出来的规律，比任何单个粒子的轨迹都可靠——粒子越多，越可靠。
\end{oneline}

\step{数学翻译器}

先给三个主角起符号。一个系统有许多微观态，第 \(i\) 个微观态出现的概率记为 \(p_i\)。所有概率加起来必须等于 \(1\)，因为“总要发生点什么”：
\[
\sum_i p_i = 1 .
\]
检查特殊情况：如果只有一个可能状态，它的概率就是 \(1\)，式子退化为 \(1=1\)，合理；如果某个 \(p_i\) 大于 \(1\)，或者全体加起来小于 \(1\)，那一定是你漏数了状态，而不是数学出了错。

任何一个随微观态变化的量 \(A\)（比如能量、正面数），它的平均值是加权求和：
\[
\langle A \rangle = \sum_i p_i A_i .
\]
检查：如果所有状态里 \(A\) 都等于同一个数 \(A_0\)，那么 \(\langle A\rangle = A_0\sum_i p_i = A_0\)，平均值就是它自己，合理。拿硬币实验做个小算例：正面数的平均值等于每个正面数乘它的排法数再相加，除以 \(1024\)。由于排法数左右对称（\(0\) 正与 \(10\) 正同为 \(1\) 种，\(4\) 正与 \(6\) 正同为 \(210\) 种），配对相消后结果正好是 \(5\)——和你在实验里用五十个数算出的平均值相去不远；抛的次数越多，越贴近 \(5\)。这也是对称性替我们省力的第一个例子：有时不用真算完所有项，看一眼分布的对称性就知道平均值落在哪。

涨落的大小用相对标准差衡量，对 \(N\) 个互不相干的个体，它约等于 \(1/\sqrt{N}\)。检查两个极端：\(N=1\) 时相对涨落约是 \(100\%\)，单次测量毫无代表性，合理；\(N\to\infty\) 时涨落趋于零，大军团给出铁打的平均值，合理——这正是开场那个胎压表的秘密，留到下一节正式结算。

然后是本章最重要的一条公式。一个与大热库接触、处于温度 \(T\) 的系统，处于能量为 \(E\) 的某个微观态的相对概率，正比于
\[
e^{-E/(k_{\mathrm B} T)},
\]
其中 \(k_{\mathrm B}\approx 1.38\times10^{-23}\,\mathrm{J/K}\) 是玻尔兹曼常数。比较两个能级，更方便的形式是比值：
\[
\frac{p_2}{p_1} = e^{-(E_2-E_1)/(k_{\mathrm B} T)} .
\]
左边回答的问题是：高能态比低能态罕见多少倍。右边的每一项都有明确含义：\(E_2-E_1\) 是要支付的“能量税”，\(k_{\mathrm B}T\) 是温度发给每个分子的“零钱尺度”——室温下约为 \(4\times10^{-21}\,\mathrm{J}\)。税相对零钱越重，付得起的状态就越少，而且少法是指数式的：税翻一倍，概率不是减半，而是按平方缩水。

立刻做四个特殊检查。第一，\(E_2=E_1\) 时比值为 \(1\)：同样贵的状态同样常见，合理。第二，\(E_2<E_1\) 时指数为正，比值大于 \(1\)：低能态更常见，方向正确。第三，\(T\to 0\) 时，只要 \(E_2>E_1\)，比值就趋于零：温度降到极低，一切都被“冻结”到能量最低的状态，这正是热力学第三定律的统计面孔。第四，\(T\to\infty\) 时比值趋于 \(1\)：温度极高时零钱无限多，各能级平起平坐。四种极限全都讲通，公式可以收下。

\begin{center}
\begin{tikzpicture}[x=1.1cm,y=0.9cm]
  \draw[-{Stealth}] (0,0) -- (6.4,0) node[right] {能量 \(E\)};
  \draw[-{Stealth}] (0,0) -- (0,3.4) node[above] {相对概率};
  \draw[thick,domain=0:6,samples=60] plot (\x,{3*exp(-\x)});
  \draw[thick,dashed,domain=0:6,samples=60] plot (\x,{3*exp(-0.35*\x)});
  \node at (4.4,0.3) {\scriptsize 低温：高能被严厉惩罚};
  \node at (5.1,1.6) {\scriptsize 高温：惩罚放宽};
  \node[left] at (0,3) {\scriptsize \(1\)};
\end{tikzpicture}
\end{center}

温度升高，曲线只是“躺平”，从不反转：高能态永远比低能态稀少，只是差距变小。这条单调性是热力学第二定律深处的一根支柱。顺带看一眼宽度：成年人身高的分布宽度大约七厘米，相对平均值不到百分之五，因为决定身高的独立因素只有那么些个；而轮胎压强的相对涨落是 \(10^{-12}\) 量级，因为背后是 \(10^{23}\) 个分子。分布的“胖瘦”，本身就是一份关于底层个体数目的情报。

有了玻尔兹曼因子，就能看懂一场贯穿全书的拔河比赛：\textbf{能量与熵的竞争}。能量想“往下沉”——低能态概率大；熵想“往多走”——微观态数目多的宏观态概率大。温度就是裁判。拿冰和水来说：冰的分子被锁在整齐的晶格里，能量低，但排列方式少；水的分子四处游走，能量高，但排列方式多到天文数字。低温时裁判偏向能量，付不起“化开”的能量税，冰赢；高温时裁判偏向熵，微观态数目的巨大优势兑现，水赢。某个确定温度上两边打平——那就是熔点。所谓自由能，就是给这场拔河记账的账本：系统自发走向的，不是能量最低的状态，也不是熵最大的状态，而是自由能最低的状态。冰在 \(0\,^\circ\mathrm{C}\) 整融化，不是分子们约好了，而是账本在这一度翻了页。

相变还带来本章最后一个新概念：\textbf{涌现}。看一块磁铁：里面有亿万个小磁针，每个只和邻居打交道——邻居朝哪，它就倾向于朝哪。高温时热运动占上风，小磁针各指各的，整块铁不显磁性；降温过了某个临界温度，邻居间的“随大流”突然滚起雪球，亿万磁针齐刷刷指向一方，磁性“冒”了出来。没有任何一根磁针下达过命令，整齐是从局部规则的无数次重复里长出来的整体性质。它像什么？像体育场里没人指挥却自发起伏的人浪。它不像什么？人浪里的人会互相观察、会犹豫会起哄，而磁针只做一件事：跟邻居对齐。正是这种“笨”规则，让物理学家能用最简单的模型算出临界温度的行为。

\begin{center}
\begin{tikzpicture}[scale=0.62]
  % 高温：无序
  \foreach \x/\r in {0/30,1/120,2/75,3/200,4/310,5/15,6/160,7/260,8/95,9/340,10/55,11/140,12/225,13/10,14/185,15/290}
    \draw[-{Stealth},thick] ({mod(\x,4)*1.1},{-int(\x/4)*1.1}) -- ++(\r:0.75);
  \node at (1.65,-4.6) {\scriptsize 高温：各指各的，整体无磁性};
  % 低温：有序
  \begin{scope}[xshift=7.2cm]
    \foreach \x in {0,...,15}
      \draw[-{Stealth},thick] ({mod(\x,4)*1.1},{-int(\x/4)*1.1}) -- ++(88:0.75);
    \node at (1.65,-4.6) {\scriptsize 低温：自发对齐，磁性涌现};
  \end{scope}
\end{tikzpicture}
\end{center}

涌现的例子一路延伸出去：交通流里没有任何一辆车想堵车，堵车却作为一种“波”在车流里传播；鸟群没有领队，却飞出整齐的队形；水分子没有“湿”这个属性，一杯水却是湿的。本章之后，从第 32 章的信息与生命，到第三册的量子统计，都会反复遇到这个主题：大军团不仅稳定，还会长出个体身上根本不存在的性质。

\begin{mathbridge}
\textbf{从分布到速率、温度和压强。}把玻尔兹曼因子用到分子的动能上，可以推出气体分子的速率分布（麦克斯韦分布）：分子速率平方的平均满足
\[
\tfrac12 m\,\overline{v^2} = \tfrac32 k_{\mathrm B} T .
\]
这给出温度的微观含义：温度是分子平均平动动能的量度——注意是“平均”和“大量”意义上的，单个分子只有动能，没有温度。代入氧气分子的质量 \(m\approx5.3\times10^{-26}\,\mathrm{kg}\) 和室温 \(T\approx300\,\mathrm{K}\)，得到方均根速率约 \(480\,\mathrm{m/s}\)，比音速还快。这解释了香水为什么一开盖就闻得到，也留下一个谜：既然分子这么快，香味为什么不是瞬间充满房间？答案是碰撞——每个分子走不了多远就撞上同伴改变方向，实际扩散是醉汉式的随机行走。这个“走不了多远”也有名字：平均自由程，即分子两次碰撞之间平均飞过的距离。常温常压下空气分子的平均自由程只有约 \(0.1\) 微米，是分子自身直径的几百倍——刚起步就撞车。于是在房间里，香味主要靠空气流动搬运，单靠扩散要等上以小时计。

再进一步，把分子撞击器壁的动量变化加起来，可以推出
\[
p = \tfrac13 n\, m\,\overline{v^2},
\]
其中 \(n\) 是单位体积的分子数。联立上面两式，就得到 \(pV = N k_{\mathrm B}T\)，正是热学章从实验总结出的理想气体定律。统计方法第一次从“数状态”里推导出了宏观定律，而不只是拟合它。
\end{mathbridge}

\begin{challenge}
\textbf{配分函数：所有可能状态的总目录。}玻尔兹曼因子只给了相对概率，要得到绝对概率，需要归一化。定义
\[
Z = \sum_i e^{-E_i/(k_{\mathrm B} T)},
\]
求和遍历系统的每一个微观态。于是 \(p_i = e^{-E_i/(k_{\mathrm B}T)}/Z\)。\(Z\) 叫配分函数，名字平淡，地位惊人：它像一本把系统所有可能状态按温度“标价”的总目录。一旦算出 \(Z\)，平均能量、熵、压强都能从它的对数求导得到，例如亥姆霍兹自由能 \(F = -k_{\mathrm B}T\ln Z\)，能量与熵的拉锯战全部浓缩在这一个函数里。熵的统计定义也在这里现身：
\[
S = k_{\mathrm B}\ln W,
\]
\(W\) 是宏观态对应的微观态数目。请严格记住这个定义：熵是“数出来的”，是微观态数目的对数，而不是一句含糊的“混乱程度”。一间整洁的病房若只有少数方式维持其整洁（\(W\) 小），它熵低；而“乱”的房间熵高，真正的原因是其微观态数目庞大，不是它“看起来乱”——看起来乱只是计数庞大在视觉上的投影。最后预告一句：当粒子数趋于无穷时，配分函数作为温度的函数可以在某个温度突然变得“不光滑”，数学上的这个尖点，就是物理上的相变。冰在 \(0\,^\circ\mathrm{C}\) 整融化，正是这个尖点的宏观表现。
\end{challenge}

\step{模型的边界}

这套“数状态”的机器威力巨大，但它的铭牌上写着四条适用条件。

第一，粒子要足够多。统计规律靠 \(1/\sqrt{N}\) 的涨落获得稳定性，\(N\) 小了，平均就失去权威。这里有一个会让直觉栽跟头的反例：我们习惯了“平均总是对的”，于是下意识认为任何一小团气体的压强、密度都和整体一样均匀。错。一立方微米（边长千分之一毫米）的空气里只有约 \(2.7\times10^{7}\) 个分子，密度涨落约 \(0.02\%\)；换成边长一百纳米的空腔，只剩几万个分子，涨落接近百分之一；到了病毒尺度的几十纳米空腔，压强会随时间明显抖动，没有任何仪表能给出稳定读数。“均匀”本身是大数效应，不是物质的天然属性。

第二，涨落并不总是看不见的，这正是直觉容易出错的另一面。1827 年，植物学家布朗在显微镜下看到水中花粉颗粒永不停歇地跳动——这是实验事实。当时有人猜是花粉“活了”，但煮死的花粉、无机尘埃照跳不误。1905 年，爱因斯坦用液体分子的无规则撞击定量解释了跳动幅度，随后佩兰的实验与预言精确吻合——“分子撞击”这个解释由此从假说升格为被验证的模型。注意三层各自的来源：跳动是看到的，分子是推断的，定量吻合是检验的。布朗运动是肉眼（借助显微镜）可见的涨落，是统计物理最早的直接证据之一。今天你在显微镜下观察一滴稀释的牛奶或墨水，同样能看到这幅景象；它提醒我们，“平静”只是尺度够大时的错觉。

第三，系统要处于或接近平衡。玻尔兹曼分布描述的是“分奖金分定了”的状态。飓风、火焰、奔腾的急流、活着的细胞，都在持续吞吐能量，是远离平衡的系统，不能直接用本章的分布去套；它们属于更难的非平衡统计物理，至今仍是前沿。就连你手边那杯正在变凉的水，严格说也只在“每一小块内部近似平衡”的意义下才能谈温度——物理学家称之为局域平衡，这是把统计工具借给非平衡世界的一座便桥。

第四，温度足够低、密度足够高时，经典数法会数错。我们把两个同种分子交换位置当作“新”的排法来计数，可在量子层面，全同粒子根本没有“谁是谁”的标记，交换不产生新状态。把不可分辨的粒子当成可分辨的来数，熵会数错——这就是历史上著名的吉布斯佯谬，它逼着物理学家在量子力学诞生之前就先修补了计数规则。修正之后的量子统计（玻色统计与费米统计）会给出与经典结果截然不同的行为，比如金属中的电子在接近绝对零度时仍保持巨大的平均能量。这是第三册的故事，这里只记住边界：本章的计数法默认粒子“可分辨、近平衡、数量大”。

还有一条边界藏在“平均”里。本章从分布推出的能量均分结果——每个运动自由度分到 \(\tfrac12 k_{\mathrm B}T\)——在分子转动、振动上也照套过，结果出了大事：照这个算法，高温物体的热辐射应在短波长处发散到无穷，史称“紫外灾难”。实验曲线却在短波端掉头向下。账本上哪里数错了？错在默认能量可以连续地、无限小份地分配。1900 年普朗克引入能量的“整份”交换，才数出和实验完全吻合的黑体辐射曲线，量子论由此开场。这是统计物理送给后续理论的最大一次“失败的礼物”：它的边界，恰好就是新物理的大门。

\step{回头解释开场现象}

现在结算胎压表。一只自行车胎内的分子数约为 \(10^{23}\) 量级，于是压强的相对涨落约 \(1/\sqrt{10^{23}}\approx3\times10^{-12}\)。仪表指针要感知这个波动，分辨率得达到小数点后十一位。指针不是不想抖，是抖动的幅度比原子还小。混乱没有消失，它只是在平均里被稀释得无影无踪——个体完全随机与整体极端稳定，不但不矛盾，反而是同一枚硬币的两面：正因为每次撞击都随机且互不相干，偏离才会按 \(\sqrt{N}\) 抵消。

再结算第二猜。一千枚硬币，正面数的标准差约为 \(\sqrt{1000}/2\approx16\) 枚，即 \(1.6\%\)。所以正面比例落在 \(49\%\) 到 \(51\%\) 之外，其实有超过四成的机会——远没有直觉以为的那么罕见。“越抛越接近一半一半”说的是相对偏差缓慢缩小，不是绝对偏差消失，更不是某个小区间内的结果变得压倒性可能。

第三猜也有了定量答案：房间空气全挤到左半边的概率约为 \(2^{-N}\)，其中 \(N\) 是 \(10^{27}\) 量级。这个概率小到连“宇宙年龄里等一次”都远远不够格。不是定律禁止它，是计数淹没了它。

顺便收获两个应用。其一，高压锅。许多烹饪和化学反应的速率正比于 \(e^{-E_{\mathrm a}/(k_{\mathrm B}T)}\)，其中 \(E_{\mathrm a}\) 是反应的“门槛能量”。取典型的门槛 \(E_{\mathrm a}\approx50\,\mathrm{kJ/mol}\)，从沸水的 \(373\,\mathrm{K}\) 升到高压锅里的 \(393\,\mathrm{K}\)，速率变为
\[
e^{\frac{E_{\mathrm a}}{R}\left(\frac1{373}-\frac1{393}\right)} \approx e^{0.82} \approx 2.3
\]
倍。区区二十度，因为撬动了指数，就把炖肉时间砍掉一半还多——工程师和厨师都在不知不觉中使用玻尔兹曼因子。其二，太阳。太阳核心温度约一千五百万度，对应的 \(k_{\mathrm B}T\) 只有约 \(1.3\,\mathrm{keV}\)，而两个质子靠近所需的库仑势垒高达 MeV 量级。平均分子根本翻不过墙，全靠分布高能尾巴里那些万里挑一的快分子，再加上量子隧穿帮忙，聚变才缓慢进行。太阳烧了近五十亿年还没烧完，恰恰是“高能态指数稀少”的功劳——分布的尾巴，决定了恒星的寿命。

其三是一个工程应用：电子器件的热噪声。电阻里的载流子也在做无规则热运动，它们数的涨落表现为电压的微小起伏——电阻两端什么都没接，也能测到随温度正比增长的噪声电压。这不是仪器粗糙，而是统计物理在电路里的直接显形。它给所有精密测量划了一条底线：麦克风的本底嘶声、射电望远镜的灵敏度、手机信号的极限，都要先和热噪声谈判。工程师的对策也很“统计”：要么降温（天文接收机常冷到液氦温度），要么多次测量取平均，用 \(1/\sqrt{N}\) 把噪声压下去。

回头看，本章的收成可以用一张清单总结：压强，是撞击次数多到涨落消失后的平均推力；温度，是平均动能这把尺子，决定能量税的税率；熵，是宏观态对应的微观态数目的对数；自由能，是能量与熵拔河时的记分牌。四个热学的核心概念，全部从“数状态”这一件事里长出来。这就是玻尔兹曼墓碑上那行公式的分量。

\step{脑洞实验与挑战题}

到这里，本章的主线已经闭合：从胎压表的静止，到微观态计数，再到玻尔兹曼因子和自由能，全部只用了一只硬币就能在家验证的逻辑。下面用一个极端尺度的思想实验把“统计词汇”的底线逼出来，再用四道挑战题检验你是否真的会用“数状态”思考。这些题目重推理、轻计算，每道都附思路提示，建议先自己争个输赢再看提示。

\begin{thought}[只剩十个分子的房间]
想象一台超级抽气机，把房间抽到只剩十个空气分子。温度计显示多少度？按定义，温度来自平均动能，你仍可以给十个分子算出某个平均值。但这个“温度”的涨落约 \(1/\sqrt{10}\approx30\%\)：刚才还是“二十度”，几个分子飞走后就变了样。再问：单个分子有温度吗？没有。它有动能，而温度是大军团平均出来的集体概念，就像“平均每个家庭 \(1.7\) 个孩子”——没有任何一个家庭真有 \(1.7\) 个孩子，但只有千家万户加在一起，这个数才有意义。继续往下抽：只剩两个分子时，“压强”变成了两次偶然撞击；只剩一个分子时，连“它服从玻尔兹曼分布”都要斟酌，因为分布说的是大军的配额，不是个体的命运。这个极端尺度的脑洞告诉我们：压强、温度、熵这些词，本质上都是统计词汇，把它们套用到少数粒子上，词本身就会失效。这也是本书最后一章要正式讨论的“涌现”：宏观性质不是微观粒子属性的简单放大，而是大军团新长出来的属性。
\end{thought}

\begin{puzzles}
\item 抛硬币连续出现十次正面。旁边的人说：“下一次该出反面了，概率会补回来。”他说得对吗？如果不对，“大数定律保证长期接近一半一半”又是什么意思？\\\textbf{思路提示：}硬币没有记忆，下一次仍然各半。大数定律不靠“补偿”实现，而靠“稀释”：已经发生的十次偏差固定不变，未来几万次里的新结果把它们淹没。均值向一半回归，靠的是分母变大，不是天平纠偏。

\item 登山到海拔约五千五百米，大气压大约只剩海平面的一半。不用查大气公式，只用本章的工具估一估：把高度差换算成一个氮气分子的重力势能差 \(mgh\)，再和 \(k_{\mathrm B}T\) 比较。\\\textbf{思路提示：}这就是玻尔兹曼因子，只不过“能量税”换成了重力势能。取 \(m\approx4.7\times10^{-26}\,\mathrm{kg}\)、\(g\approx10\,\mathrm{m/s^2}\)、\(h\approx5.5\times10^3\,\mathrm{m}\)、\(T\approx270\,\mathrm{K}\)，算出 \(mgh/(k_{\mathrm B}T)\approx0.7\)，于是 \(e^{-0.7}\approx0.5\)。大气本身就是一座按玻尔兹曼分布堆起来的分子塔。

\item 拿一个水分子问它：“你是冰还是蒸汽？”这个问题有意义吗？为什么水偏偏在 \(100\,^\circ\mathrm{C}\) 整沸腾，而不是从 \(95\,^\circ\mathrm{C}\) 开始慢慢“有点沸腾”？\\\textbf{思路提示：}冰和蒸汽都不是单个分子的属性，而是分子之间排列关系的属性。相变是能量与熵竞争、再加上分子间相互作用的集体事件；粒子数趋于无穷时，配分函数在特定温度出现数学上的不光滑，于是宏观上呈现一条锋利的界线。“一个水分子”既不在界线的这一边，也不在那一边——界线属于大军团。

\item 两杯各 \(50\,^\circ\mathrm{C}\) 的水倒在一起，为什么不是 \(100\,^\circ\mathrm{C}\)？反过来，半杯 \(50\,^\circ\mathrm{C}\) 的水和整杯 \(50\,^\circ\mathrm{C}\) 的水，哪个“更热”？\\\textbf{思路提示：}温度是平均动能，是强度量：把两份相同的系统合并，分子数翻倍，平均值不变。会翻倍的是能量和微观态数目这类广延量。分清“随人数翻倍的量”和“与人数无关的量”，是统计物理送给日常语言的最好礼物。
\end{puzzles}
